\section{Discussion}
\label{sec:discussion}

\subsection{Both agents found the bottleneck; only one kept going}

Neither agent was told where to look, and both went to crossing minimization.
Working in isolation, both arrived at the same transformation: fusing the two
symmetric crossing-count passes. Locating the expensive phase is evidently not
the hard part.

What separates them is how far they went next. Codex stopped at the phase it
identified. Claude continued into network simplex, which serves both ranking and
coordinate assignment, and that accounts for its entire margin --- the same
pattern on three workloads with very different composition, not an inference
from one measurement.

\subsection{Why breadth mattered: the ceiling moves with shape}

Amdahl's law~\cite{amdahl1967validity} bounds what optimizing one phase can
yield, and that bound depends on the phase's share of runtime --- which
Table~\ref{tab:composition} shows is a property of the graph's shape.

\begin{table}[htbp]
\centering
\caption{What optimizing crossing minimization alone can achieve, against what
each agent measured. \textbf{Ceiling:} baseline wall clock divided by (wall
clock minus that phase's measured cost), i.e.\ the phase driven to zero.
\textbf{Claude and Codex:} the end-to-end speedups from
Table~\ref{tab:endtoend}, carrying the same dagger convention.}
\label{tab:ceiling}
\begin{tabular}{lrrr}
\toprule
\textbf{Graph} & \textbf{Ceiling, crossing min.\ only} & \textbf{Claude} & \textbf{Codex} \\
\midrule
sparse-deep   & 1.11$\times$ & \textbf{1.33$\times$} & 0.99$\times$$^\dagger$ \\
default       & 1.69$\times$ & 1.50$\times$ & 1.16$\times$ \\
dense-shallow & 2.58$\times$ & 1.84$\times$ & 1.29$\times$ \\
\bottomrule
\end{tabular}
\end{table}

Codex sits below the ceiling on all three shapes, as an agent that optimized
only that phase must. On sparse-deep the ceiling is $1.11\times$, so even a
perfect crossing-minimization optimization would be nearly invisible end-to-end.
Codex's $0.99\times$ is not a failed optimization --- it works, at $1.31\times$
on the phase --- but a phase worth $11.5\%$ of the workload.

Claude \emph{exceeds} that ceiling, $1.33\times$ against $1.11\times$, which is
only possible by improving other phases. That is the breadth argument at its
sharpest: on that workload no amount of further work on the phase both agents
found could have matched it.

The ceiling an optimization can reach is a property of the input, not the
program, and a single benchmark graph cannot show that.

\subsection{Cost}

\begin{table}[H]
\centering
\caption{Claude run cost, as reported by its session log.}
\label{tab:cost-claude}
\begin{tabular}{lr}
\toprule
\textbf{Metric} & \textbf{Value} \\
\midrule
Model                 & \texttt{claude-opus-5} \\
Working time          & 4\,h\,27\,m\,3\,s \\
\midrule
Input tokens          & 730 \\
Output tokens         & 620,448 \\
Cache creation tokens & 1,270,669 \\
Cache read tokens     & 89,851,473 \\
\midrule
Total tokens          & 91,743,320 \\
\bottomrule
\end{tabular}
\end{table}

\begin{table}[H]
\centering
\caption{Codex run cost, as reported by its session log. Cached input is a
subset of input; reasoning output is a subset of output.}
\label{tab:cost-codex}
\begin{tabular}{lr}
\toprule
\textbf{Metric} & \textbf{Value} \\
\midrule
Model                     & \texttt{gpt-5.6-terra} \\
Working time              & 23\,m\,42\,s \\
\midrule
Input tokens              & 5,114,035 \\
\quad Cached input tokens & 4,985,600 \\
Output tokens             & 20,826 \\
\quad Reasoning tokens    & 4,555 \\
Cache-write tokens        & 0 \\
\midrule
Total tokens              & 5,134,861 \\
\bottomrule
\end{tabular}
\end{table}
